\documentclass[11pt]{article}
\usepackage[margin=0.8in]{geometry}
\usepackage{graphicx}
\usepackage{amsmath}
\usepackage{hyperref}
\title{n=47 random00: the degree-four closure equals the accepted quotient kernel}
\author{Exact finite-instance evidence record}
\date{2 November 2026}
\begin{document}
\maketitle
\noindent The independent implementation B replay returned degree-four Boolean closure rank $151{,}857$ in $B_4$ of dimension $213{,}053$, with $1\notin W_4(F)$. The original independent D4 checker accepted a surjective quotient with $61{,}196$ coordinates and proved $W_4(F)\subseteq U_4$. Therefore
\[
\dim U_4=213{,}053-61{,}196=151{,}857=\dim W_4(F),
\qquad W_4(F)=U_4.
\]
The conclusion applies to the single frozen n=47 random00 Boolean presentation. The replay reported $12{,}535$ seconds under a 16 GiB child address-space cap; no comparative timing claim is made.
\begin{center}
\includegraphics[width=0.96\textwidth]{dimension-bridge.pdf}
\end{center}
\noindent The frozen plan pins the field, curve, target, subspace, accepted seed hash, checker transcript and both source hashes. The archived output has status \texttt{rank\_match}, return code 0, no timeout, and empty stderr. \texttt{audit.py} verifies every archived output digest and, with the local seed, the 1.63 GB seed hash and header. The audit does not rerun the exact closure.

\medskip\noindent This closes the D4 seed-equality obligation for this input. The earlier D5 producer reported a killed constant but retained no independently checkable terminal relation witness. Hence this report establishes $U_4=W_4$ and leaves the D5 membership/upper-bound question open. A checked D5 primal derivation or independent full relation replay is the next proof step.

\medskip\noindent Evidence: the frozen plan, archived result and checker transcript, evidence manifest, and audit script in this report directory. The D4 seed SHA-256 begins \texttt{3bdad807113e43ec}; the complete digest is in the plan and manifest.
\end{document}
